The wall of the anisotropic case with the helix of the genet23 test case
(\texttt{tests/cases/chamber\_cylinder.json}): $+60^\circ$ at the endocardium to $-60^\circ$
at the epicardium. Helical fibers load every cross-section of the uniform 0D state with
the azimuthal shear $\sigma_{\theta z}(R)$: positive near the endocardium, negative near the
epicardium, zero net torque (which sets $\beta$), up to 18\% of $\sigma_{\theta\theta}$ at
120~mmHg. The 0D block has no ends, so this costs nothing; a 3D end that only keeps
$u_z = 0$ carries no $\sigma_{\theta z}$, and the 3D tube twists and expands differently
in end layers (Figures in the appendix). The 3D model is therefore the full cylinder
with free ends, and only its mid-length section is compared.

\paragraph{Tube length.} All helical comparisons use a tube twice as long as that of the
test case ($2L_0 = 114$~mm, 8.2 wall thicknesses; the 0D state does not depend on the
length). In the original tube the end layers reach the middle: at nr = 4 its mid-length
radial displacement differs from that of the $2L_0$ tube by $3.4\cdot10^{-3}$ at the
endocardium and its twist rate by 14\%, and both stop converging to the 0D block. A tube
four times as long gives the mid-length state of the $2L_0$ tube to $10^{-4}$ in the radial
displacement at nr = 2 (and to $2\cdot10^{-5}$ for the uniform $30^\circ$ helix), so the
middle of the $2L_0$ tube is free of the end effect. The dashpots that hold the rigid-body
modes do not matter either: ten times or a tenth of their damping changes the mid-length
radial displacement by $1.5\cdot10^{-5}$ and $1.5\cdot10^{-6}$, and dashpots on the
epicardium instead of the ends give the same state to $10^{-5}$.

\paragraph{Results.} At the mid-length section of the $2L_0$ tube every quantity converges
to the 0D block (Table~\ref{tab:passive_helical_L2-kin}): the volume at order 3.8 and 2.8 to
$1.7\cdot10^{-4}$ at nr = 4 (38016 elements, on the cluster), the radial displacement likewise
to $8.7\cdot10^{-4}$, the wall thickness at 2.1--2.8, the transmural rotation at 2.6--3.3, the
strains and stresses at 1.6--2.2. The twist rate converges late, 15\%, 8\% and 1.6\% at
nr = 1, 2, 4 (3D/0D = 1.016 at nr = 4): the $\pm60^\circ$ helix twists the tube by only
$-0.17^\circ$ over its length, a small difference of the torques of its two fiber layers
that amplifies the stress error. The two helices below, with large twists, show the
torsion of the 0D block without this cancellation.
